%%%%%%%%%%%%%%%%%%%%%%%%%%%%%%%%%%%%%%%%%%%%%%%%%%%%
\documentclass[fleqn,10pt,twocolumn]{AROB}
\usepackage{amssymb}
\usepackage{amsmath}

\title{Streaming Apriori-Window: Mining Dense Patterns\\
       from Streaming Data for Autonomous Agent Monitoring}

\author{Kanata Takayasu${}^{1\dagger}$, Taihei Takahashi${}^{1}$,
        Satoshi Suga${}^{2}$, and Satoshi Kurihara${}^{1}$}
% The dagger symbol indicates the presenter.
\speaker{Kanata Takayasu}

\affils{${}^{1}$Keio University, Yokohama, Japan\\
${}^{2}$Kansai University, Osaka, Japan\\
}

\abstract{%
Autonomous AI systems and robotic agents generate streams of sensor events
whose co-occurrence patterns reveal behavioral regularities.
Detecting \textit{when} these patterns densely co-occur requires a
streaming-capable interval mining framework with controllable output latency.
We propose Streaming Apriori-Window, parameterized by a scan frequency $n$,
which identifies all time intervals where itemset frequency in a
fixed-width sliding window meets a threshold.
A confirmed output guarantee follows from the monotone scan property.
Two speedup techniques --- stride-based skipping and Apriori-style
candidate pruning --- achieve practical processing speed.
Experiments show that $n = W$ achieves a 15.5$\times$ speedup over $n = 1$
with a 24-hour scan interval, and that results on CASAS smart home sensor
data are interpretable and domain-consistent.%
}

\keywords{%
data stream mining, pattern recognition, sensor data analysis,
behavior analysis, intelligent monitoring%
}

\begin{document}

\maketitle

%-----------------------------------------------------------------------

%%%%%%%%%%%%%%%%%%%%%%%%%%%%%%%%%%%%%%%%%%%%%%%%%%%%
\section{Introduction}

Autonomous AI systems and robotic agents generate streams of sensor events
that encode behavioral patterns.
Understanding \textit{when} these patterns densely co-occur --- not merely whether
they exist globally --- is essential for real-time monitoring, anomaly detection,
and adaptive control in intelligent environments.
In multi-agent systems and sensor-equipped robotic platforms,
detecting that a specific set of events repeatedly clusters within particular
time intervals can reveal coordinated behaviors, routine structures,
or emerging anomalies that demand timely response.
The approach also has broad applicability across domains such as retail and epidemiology.

We propose \textit{Streaming Apriori-Window}, a dense interval mining
algorithm for event streams parameterized by a scan frequency $n$.
The algorithm identifies all time intervals in which itemset occurrence
frequency within a fixed-width sliding window meets a threshold.
Two speedup techniques --- stride-based skipping and Apriori-style
candidate pruning --- are incorporated to achieve practical processing speed.

The main contributions of this paper are as follows:
\begin{itemize}
  \item We formalize the streaming dense interval mining problem with
        an explicit scan frequency parameter $n$ that controls the tradeoff
        between output latency and processing cost.
  \item We prove a confirmed output guarantee: intervals ending before
        the current scan position are unaffected by future events.
  \item We evaluate the algorithm on synthetic and real-world CASAS
        smart home sensor data, demonstrating practical processing speed
        and interpretable outputs.
\end{itemize}


%%%%%%%%%%%%%%%%%%%%%%%%%%%%%%%%%%%%%%%%%%%%%%%%%%%%
\section{Related Work}

\subsection{Frequent Itemset Mining}

Apriori~\cite{agrawal1994fast} and FP-Growth~\cite{han2000mining}
are foundational methods for mining frequent itemsets.
ECLAT~\cite{zaki1997new} offers a vertical data format alternative.
A comprehensive survey of FIM is given in~\cite{fournier2017survey}.
All of these methods evaluate support over the entire dataset and
cannot detect itemsets that occur at high frequency only within
specific time intervals.

\subsection{Dense Interval Mining}

Methods that output patterns together with their dense intervals
fall into two categories.
\textit{Occurrence-gap-based} methods include
LPFIM~\cite{yin2014lpfim}, LPPM~\cite{saleh2011lppm}, and
RPM~\cite{kiran2017rpm}.
LPFIM partitions intervals when occurrence gaps reach a threshold;
LPPM introduces cumulative spillover for noise-robust detection;
RPM groups consecutive occurrences with gaps within a maximum period.
These methods indirectly estimate density through gap constraints
or periodicity thresholds, making parameter tuning difficult and
results inexact with respect to the direct sliding-window frequency.
\textit{Direct} approaches include GLFMiner~\cite{fournierviger2021glf},
which mines patterns at fixed temporal granularities and thus cannot
identify intervals of arbitrary width, and SIM~\cite{kiran2015sim},
which returns only the single optimal interval per itemset.
None of these methods support streaming or provide a mechanism
to control scan frequency.

\subsection{Stream Mining}

Sliding-window-based stream mining has been studied in
incremental FIM~\cite{lee2001sliding,leung2006dstree},
where windows limit the temporal scope of support computation.
However, these methods still report global support within the window
and do not enumerate all dense sub-intervals.



%%%%%%%%%%%%%%%%%%%%%%%%%%%%%%%%%%%%%%%%%%%%%%%%%%%%
\section{Problem Formulation}

Let $\mathcal{I}$ be a finite set of items.
An \textbf{event stream} is a sequence of transactions
$T_0, T_1, T_2, \ldots$, where each $T_t \subseteq \mathcal{I}$.
We associate each transaction with a discrete timestamp $t \in \mathbb{Z}_{\geq 0}$.

\subsection{Dense Interval}

Given a window width $W \in \mathbb{Z}_{>0}$ and a threshold $\theta \in \mathbb{Z}_{>0}$,
the \textbf{occurrence frequency} of an itemset $X \subseteq \mathcal{I}$
in window $[l, l+W)$ is
\[
  \mathrm{freq}(X, l) = |\{t \in [l, l+W) : X \subseteq T_t\}|.
\]
A window $[l, l+W)$ is \textbf{dense} for $X$ if $\mathrm{freq}(X, l) \geq \theta$.
A \textbf{dense interval} for $X$ is a maximal contiguous range
$[l_s, l_e]$ such that $[l, l+W)$ is dense for $X$ for all $l \in [l_s, l_e]$.

\subsection{Streaming Dense Interval Mining}

Given an event stream $T_0, T_1, \ldots$, parameters $W$, $\theta$, and
a scan frequency $n \in \mathbb{Z}_{>0}$, the task is to enumerate,
at each scan step $s = 0, n, 2n, \ldots$, all dense intervals
$(X, [l_s, l_e])$ whose right endpoint satisfies $l_e + W \leq s$,
without revisiting previously output intervals.
We call such output \textbf{confirmed}: because the sliding window
advances monotonically, intervals ending before position $s$
cannot change with future transactions.

The parameter $n$ controls the tradeoff:
$n = 1$ yields minimal latency at maximum cost;
$n = T$ (full stream length) reduces to batch processing.


%%%%%%%%%%%%%%%%%%%%%%%%%%%%%%%%%%%%%%%%%%%%%%%%%%%%
\section{Algorithm}

\subsection{Monotone Scan}

For a fixed itemset $X$, let $\mathbf{p} = (p_0, p_1, \ldots)$ be the sorted
list of timestamps at which $X$ occurs.
We scan the left endpoint $l$ from $0$ to the current position,
maintaining a pointer into $\mathbf{p}$.

\textbf{Key property}: as $l$ increases by one, the window $[l, l+W)$
loses $T_{l-1}$ and may gain $T_{l+W-1}$.
Consequently, $\mathrm{freq}(X, l)$ changes by at most $\pm 1$ per step,
and the left endpoint $l$ is monotonically non-decreasing.
This means any confirmed output is never invalidated.

\subsection{Stride-Based Skipping}

Na\"ively incrementing $l$ by 1 each step is wasteful.
We exploit two cases:

\textbf{Dense phase}: if $\mathrm{freq}(X, l) \geq \theta$,
the next non-dense position is at least $l + \mathrm{surplus}$, where
$\mathrm{surplus} = \mathrm{freq}(X, l) - \theta + 1$.
We skip ahead by $\mathrm{surplus}$ steps.

\textbf{Non-dense phase}: if $\mathrm{freq}(X, l) < \theta$,
the next potentially dense position requires
$\theta - \mathrm{freq}(X, l)$ additional occurrences.
We advance $l$ to the position of the $(\theta - \mathrm{deficit})$-th
next occurrence, skipping the deficit gap.

\subsection{Apriori-Style Candidate Pruning}

For itemset $X$ of size $k$, if any $(k{-}1)$-subset $Y \subset X$
has no dense interval intersecting $[l, l+W)$, then $X$ cannot be
dense in $[l, l+W)$ either (by the anti-monotonicity of support).
We precompute, for each $(k{-}1)$-subset, the set of $l$-values that
are dense, and restrict candidate scan ranges for $X$ to the
intersection of its subsets' dense ranges.

\subsection{n-Step Streaming}

In the target deployment, transactions arrive continuously from a
sensor stream or autonomous agent log and are appended to a database.
The system waits until $n$ new transactions have accumulated since the
last scan, then triggers a mining step.
The parameter $n$ directly controls output latency: smaller $n$ yields
more responsive output at higher computational cost; larger $n$ batches
more events per scan, reducing cost.
Setting $n = T$ (full stream length) recovers batch processing as a
special case.

Each scan step proceeds as follows:

\begin{enumerate}
  \item When $n$ new transactions have arrived (current position $s$),
        retrieve the buffer $[s - W + 1,\; s + n)$, retaining the
        $W{-}1$ most recent transactions from the previous step to
        prevent boundary loss.
  \item Run the monotone scan with stride-based skipping and
        Apriori-style pruning over this buffer.
  \item Emit all intervals $(X, [l_s, l_e])$ with $l_e + W \leq s$
        not yet output: these are \textbf{confirmed} and will not
        change as future transactions arrive.
  \item Advance to the next accumulation of $n$ transactions and repeat.
\end{enumerate}


%%%%%%%%%%%%%%%%%%%%%%%%%%%%%%%%%%%%%%%%%%%%%%%%%%%%
\section{Experiments}

All experiments were conducted on a single machine
(Apple M3 Pro, 18 GB RAM) using Rust~1.82 release builds.
The CASAS dataset~\cite{casas2013} records timestamped ON/OFF events
from motion and door sensors in a smart home,
discretized into 1-hour slots.

\textbf{Benchmark methodology.}
In production, computation time and transaction arrival time are
interleaved: the system waits for $n$ transactions, then computes.
To measure \textit{pure computation time} independently of arrival
latency, we pre-load the full dataset and feed it to the system
$n$ transactions at a time, triggering each scan immediately after
the previous one completes.
This isolates algorithmic cost from network or sensor delays and
yields a fair comparison across different values of $n$.

\subsection{EX-1: Scan Frequency vs.\ Processing Speed}

We varied $n \in \{1, 6, 24, 96, T\}$ on CASAS hh113
($T = 12{,}226$, $W = 24$, $\theta = 10$, $k = 4$)
and measured total processing time.

\begin{table}[t]
\caption{Processing speed vs.\ scan frequency ($n$) on CASAS hh113.}
\begin{center}
\begin{tabular}{|c|r|r|}
\hline
$n$ & Time (ms) & Speedup vs.\ $n{=}1$ \\\hline
1   & 1,068     & 1.0$\times$ \\
6   &   216     & 4.9$\times$ \\
24  &    69     & 15.5$\times$ \\
96  &    27     & 39.6$\times$ \\
$T$ (batch) & 7.5 & 142.4$\times$ \\\hline
\end{tabular}
\end{center}
\end{table}

$n = W = 24$ achieves a 15.5$\times$ speedup over $n = 1$
while scanning every 24 hours, a practically useful frequency.
The improvement saturates after $n = W$,
suggesting $n = W$ as the recommended default.

\subsection{EX-2: Ablation Study}

We measured the contribution of each speedup technique
on synthetic data ($T = 12{,}226$, $W = 24$, $\theta = 10$, $k = 4$).

\begin{table}[t]
\caption{Ablation: processing time [ms] for four algorithm variants.}
\begin{center}
\begin{tabular}{|l|c|c|r|}
\hline
Variant & Stride & Pruning & Time (ms) \\\hline
Naive          &        &         & TODO \\
Stride-only    & \checkmark &    & TODO \\
Pruning-only   &        & \checkmark & TODO \\
Full (proposed)& \checkmark & \checkmark & TODO \\\hline
\end{tabular}
\end{center}
\end{table}

% TODO: fill in ablation results after running experiments

\subsection{EX-3: Application to CASAS Sensor Data}

We applied the algorithm ($n = 24$, batch mode) to two CASAS datasets.
CASAS records multi-sensor co-activation streams structurally analogous
to autonomous robotic behavioral logs.

\begin{table}[t]
\caption{CASAS dataset statistics and detected dense intervals.}
\begin{center}
\begin{tabular}{|l|r|r|r|}
\hline
Dataset & $T$ & $W$ & Patterns \\\hline
hh113 & 12,226 & 24 & 57 \\
mv001 & 34,759 & 24 & 53 \\\hline
\end{tabular}
\end{center}
\end{table}

% TODO: add representative pattern examples and interpretation


%%%%%%%%%%%%%%%%%%%%%%%%%%%%%%%%%%%%%%%%%%%%%%%%%%%%
\section{Conclusion}

We proposed Streaming Apriori-Window, a dense interval mining algorithm
for event streams with an explicit scan frequency parameter $n$.
The confirmed output guarantee follows from the monotone sliding-window
property, ensuring that output intervals are never invalidated by
future events.
Two speedup techniques --- stride-based skipping and Apriori-style
candidate pruning --- reduce processing cost significantly,
with $n = W$ achieving a favorable tradeoff between latency and speed.
Experiments on CASAS smart home sensor data demonstrate that the
algorithm produces interpretable, domain-consistent outputs,
confirming its practical utility for autonomous agent monitoring.

Future work includes applying the framework to behavioral logs from
physical robotic systems, extending it to multi-dimensional event spaces,
and investigating adaptive scan-frequency mechanisms based on
stream characteristics.


%-----------------------------------------------------------------------
%%%%%%%%%%%%%%%%% BIBLIOGRAPHY IN THE LaTeX file !!!!! %%%%%%%%%%%%%%%%%%%%%%

\begin{thebibliography}{99}

\bibitem{agrawal1994fast}
R. Agrawal and R. Srikant,
``Fast algorithms for mining association rules,''
{\it Proc. VLDB}, pp.~487--499, 1994.

\bibitem{zaki1997new}
M. J. Zaki, S. Parthasarathy, M. Ogihara, and W. Li,
``New algorithms for fast discovery of association rules,''
{\it Proc. KDD}, pp.~283--296, 1997.

\bibitem{han2000mining}
J. Han, J. Pei, and Y. Yin,
``Mining frequent patterns without candidate generation,''
{\it ACM SIGMOD Record}, vol.~29, no.~2, pp.~1--12, 2000.

\bibitem{fournier2017survey}
P. Fournier-Viger, J. C.-W. Lin, B. Vo, T. T. Chi, J. Zhang, and H. B. Le,
``A survey of itemset mining,''
{\it WIREs Data Mining Knowl. Discov.}, vol.~7, no.~4, e1207, 2017.

\bibitem{yin2014lpfim}
K.-C. Yin, Y.-L. Hsieh, D.-L. Yang, and M.-C. Hung,
``Association rule mining considering local frequent patterns with temporal intervals,''
{\it Applied Mathematics \& Information Sciences}, vol.~8, no.~4,
pp.~1879--1890, 2014.

\bibitem{saleh2011lppm}
B. Saleh and F. Masseglia,
``Discovering frequent behaviors: Time is an essential element of the context,''
{\it Knowledge and Information Systems}, vol.~28, no.~2, pp.~311--331, 2011.

\bibitem{kiran2017rpm}
R. U. Kiran, J. Venkatesh, M. Toyoda, M. Kitsuregawa, and P. K. Reddy,
``Discovering partial periodic-frequent patterns in a transactional database,''
{\it J. Systems and Software}, vol.~125, pp.~170--182, 2017.

\bibitem{fournierviger2021glf}
P. Fournier-Viger {\it et al.},
``Mining local periodic patterns in a discrete sequence,''
{\it Information Sciences}, vol.~544, pp.~519--548, 2021.

\bibitem{kiran2015sim}
R. U. Kiran, H. Shang, M. Toyoda, and M. Kitsuregawa,
``Discovering recurring patterns in time series,''
{\it Proc. EDBT}, pp.~97--108, 2015.

\bibitem{lee2001sliding}
C.-H. Lee, C.-R. Lin, and M.-S. Chen,
``Sliding-window filtering: An efficient algorithm for incremental mining,''
{\it Proc. CIKM}, pp.~263--270, 2001.

\bibitem{leung2006dstree}
C. K.-S. Leung and Q. I. Khan,
``DSTree: A tree structure for the mining of frequent sets from data streams,''
{\it Proc. ICDM}, pp.~928--932, 2006.

\bibitem{casas2013}
D. Cook,
``CASAS: A smart home in a box,''
{\it IEEE Computer}, vol.~46, no.~7, pp.~62--69, 2013.

\end{thebibliography}

\end{document}
